%%
%% Copyright 2019-2020 Elsevier Ltd
%%
%% This file is part of the 'CAS Bundle'.
%% --------------------------------------
%%
%% It may be distributed under the conditions of the LaTeX Project Public
%% License, either version 1.2 of this license or (at your option) any
%% later version.  The latest version of this license is in
%%    http://www.latex-project.org/lppl.txt
%% and version 1.2 or later is part of all distributions of LaTeX
%% version 1999/12/01 or later.
%%
%% The list of all files belonging to the 'CAS Bundle' is
%% given in the file `manifest.txt'.
%%
%% Template article for cas-sc documentclass for
%% double column output.

%\documentclass[a4paper,fleqn,longmktitle]{cas-sc}
\documentclass[a4paper,fleqn]{cas-sc}
\usepackage{njfutitlepage}
% \usepackage[numbers]{natbib}
\usepackage[authoryear]{natbib}
% \usepackage[authoryear,longnamesfirst]{natbib}
\usepackage{csquotes}
\usepackage[UTF8]{ctex}
\usepackage[ruled,vlined]{algorithm2e}
\usepackage{etoolbox}
\usepackage{placeins}

% \newcounter{stage}
% \setcounter{stage}{0}
% \newcommand{\Stage}[1]{%
%   \stepcounter{stage}%
%   \vspace{0.3em}%
%   \textbf{\textit{\footnotesize\arabic{stage}. #1:\\}}%
%   \vspace{0.2em}%
% }
% \AtEndEnvironment{algorithm}{\setcounter{stage}{0}}

\newcommand{\refSec}[1]{\hyperref[#1]{Section~\ref*{#1}}}
\newcommand{\refFig}[1]{\hyperref[#1]{Fig.~\ref*{#1}}}
\newcommand{\refApp}[1]{\hyperref[#1]{Appendix~\ref*{#1}}}
\renewcommand{\figurename}{Fig.}

\newcommand{\FeHALSRepo}{\href{https://github.com/WZYivan/FeHALS}{FeHALS github Main Repo}}

%%%Author definitions
\def\tsc#1{\csdef{#1}{\textsc{\lowercase{#1}}\xspace}}
\tsc{WGM}
\tsc{QE}
\tsc{EP}
\tsc{PMS}
\tsc{BEC}
\tsc{DE}
%%%

\begin{document}
\let\WriteBookmarks\relax
\def\floatpagepagefraction{1}
\def\textpagefraction{.001}

\shortauthors{7组}
\shorttitle{}
\title [mode = title]{FeHALS：基于Web的3D可视化航路规划与激光仿真系统}

\Course{\enquote{智能+}GIS设计与开发}
\Major{测绘工程}
\Grade{2023级}
\Group{7 组}

\Member[2310604116]{王 孜悠}
\Member[2310604109]{梅 康凯}
\Member[2310604115]{唐 如意}
\Member[2310604122]{杨 杰瑞}
\Member[2310604117]{温 璞}
\Member[2110604112]{刘 继明}
\Member[2110604106]{李 英杰}
\Member[2310604103]{杜 仲宇}
\Member[2310604113]{施 祺}
\Member[2310604118]{向 宇鸣}

\Teacher{隋 铭明}
\Teacher{那 嘉明}
\Teacher{陈 动}
\Teacher{朱 杰}

\author[1]{王 孜悠}[style=chinese, orcid=2310604116]
\ead{WZYivan}
\credit{文档编写, 思路构建, 代码实现, 功能验证, 项目管理, 代码审查, 图表绘制}

\author[1]{梅 康凯}[style=chinese, orcid=2310604109]
\ead{zhangqilinghy}
\credit{代码实现, 系统适配}

\author[1]{唐 如意}[style=chinese, orcid=2310604115]
\ead{Try6743}
\credit{代码实现}

\author[1]{杨 杰瑞}[style=chinese, orcid=2310604122]
\ead{19951391030yjr-sudo}
\credit{文档编写, 图表绘制}

\author[1]{杜 仲宇}[style=chinese, orcid=2310604103]
\ead{2309175724@qq.com}
\credit{文档编写, 图表绘制}

\author[1]{李 英杰}[style=chinese, orcid=2110604106]
\ead{1589428766@qq.com}
\credit{代码实现, 系统适配}

\author[1]{施 祺}[style=chinese, orcid=2310604113]
\ead{69901001+nidie1@users.noreply.github.com}
\credit{代码实现}

\author[1]{温 璞}[style=chinese, orcid=2310604117]
\credit{暂无贡献}

\author[1]{刘 继明}[style=chinese, orcid=2110604112]
\credit{暂无贡献}

\author[1]{向 宇鸣}[style=chinese, orcid=2310604118]
\credit{暂无贡献}

\author[1]{隋 铭明}[style=chinese]
\fnmark[1]
\fntext[1]{指导老师}
\credit{监督指导}

\author[1]{那 嘉明}[style=chinese]
\fnmark[1]
\credit{监督指导}

\author[1]{陈 动}[style=chinese]
\fnmark[1]
\credit{监督指导}

\author[1]{朱 杰}[style=chinese]
\fnmark[1]
\credit{监督指导}

\affiliation[1]{organization={College of Civil Engineering, Nanjing Forestry University},
    addressline={Longpan Rd.},
    city={Nanjing},
    postcode={210037},
    country={China}}

\begin{abstract}
    三维激光扫描仿真技术在地理信息、遥感与测绘领域的应用日益广泛。HELIOS++作为一款开源的激光雷达仿真框架，能够高保真地模拟机载、地面与无人机等平台的扫描过程，但其原生交互方式依赖命令行与XML配置文件，缺乏直观的三维可视化航路规划能力，限制了其在复杂场景下的高效使用。本文提出并实现了一个基于Web的FeHALS（Frontend of HELIOS++ for Airborne Laser Scanning）系统。系统采用浏览器/服务器（B/S）架构，前端基于Vue 3与Three.js构建三维交互场景，后端基于FastAPI提供仿真服务并封装HELIOS++引擎调用。系统支持导入多种三维模型，用户可在浏览器三维场景中通过鼠标交互设定航点并自动生成符合HELIOS++格式的航迹文件；同时提供图形化界面配置扫描平台与传感器参数，一键调用仿真引擎并通过WebSocket实时推送运行日志；仿真生成的点云可自动加载至三维场景，并支持按高度、强度着色以及点尺寸与透明度的调节。在项目管理方面，本文作者作为仓库唯一的CODEOWNER，全面负责系统架构设计、技术路线决策、核心代码实现、文档编写与团队协作的统筹管理，通过GitHub的Pull Request与分支保护机制确保代码质量与协作效率。该系统有效降低了HELIOS++的使用门槛，为航摄方案设计、传感器参数优化及点云数据模拟提供了便捷、高效的可视化工具。本文详细阐述了系统的架构设计、核心功能模块与技术实现路径，并通过典型场景验证了系统的可行性与实用性。源代码已公开至\url{https://github.com/WZYivan/FeHALS}。
\end{abstract}

\begin{graphicalabstract}
    \includegraphics[width=\textwidth]{figs/graphicalabstract.pdf}
\end{graphicalabstract}

% Research highlights
\begin{highlights}
    \item 设计并实现了基于Web的HELIOS++仿真参数图形化配置界面，自动生成XML配置文件与航迹文件。
    \item 构建了从场景构建、航路设计、仿真执行到结果渲染的完整闭环工作流，实现了激光仿真全流程的集成化与可视化。
    \item 系统基于开源技术栈（Vue 3 + Three.js + FastAPI + HELIOS++），前后端解耦，具有良好的可扩展性与跨平台访问能力。
    \item 使用git+github的本地+云端组合统筹开发流程，深度落实了现代化软件开发协作机制。
\end{highlights}

% Keywords
\begin{keywords}
    ALS \sep 航路规划 \sep 三维可视化 \sep B/S架构 \sep HELIOS++
\end{keywords}

\MakeNJFUTitlePage[UseDocTitle]
\maketitle
\tableofcontents
\hrulefill

\section{引言}
激光雷达（Light Detection and Ranging，LiDAR）通过发射激光脉冲并测量其回波时间与强度，能够快速获取地物表面的高精度三维点云，已成为航空摄影测量、林业资源调查、城市三维建模与地形测绘等领域的核心数据获取手段之一 \citep{laserRemoteSensing}。机载激光扫描（Airborne Laser Scanning，ALS）将激光扫描仪搭载于飞行平台，沿预定航线对地面进行扫描，其点云质量与航高、航速、扫描频率、脉冲频率及航线布设等参数密切相关。然而，实际航摄作业受空域、成本与环境等因素制约，方案设计与参数验证往往缺乏充分的试验条件，因此通过计算机仿真对扫描过程进行预演、评估与优化具有重要的工程意义。

HELIOS++ \citep{heliosPlusPlus} 是一款开源的激光雷达仿真框架，基于光线追踪算法能够高保真地模拟机载、地面、车载与无人机等平台的扫描过程，并支持全波形输出。但其原生使用方式依赖命令行与手写XML配置文件，用户需要熟练掌握场景（scene）、扫描任务（survey）、平台（platform）与扫描器（scanner）等配置文件的语法，且无法直观地在三维场景中规划航线，学习与使用门槛较高。

HELIOS++官方团队已提供了一个基于QGIS的图形化插件AEOS \citep{aeos}，将HELIOS++的仿真配置与执行功能集成到QGIS桌面环境中。AEOS作为QGIS插件，充分利用了QGIS的三维地图渲染与地理空间数据管理能力，适合需要与GIS工作流深度结合的用户。然而，AEOS的使用依赖于QGIS桌面环境的完整安装与配置，对于仅需快速进行航路规划与仿真参数调试的用户而言，启动和配置成本较高，且无法在浏览器或轻量级环境中直接使用。

针对上述问题，本文设计并实现了FeHALS（Frontend of HELIOS++ for Airborne Laser Scanning），一个基于Web的三维可视化航路规划与激光仿真系统，其Web前端如\refFig{fig:fehals_overview}所示。与AEOS的QGIS集成路线不同，FeHALS采用浏览器/服务器（B/S）架构：前端基于Vue 3 \citep{vuejs} 与Three.js \citep{threejs} 构建三维交互场景，提供模型加载、交互式航点编辑、参数配置与点云渲染等功能；后端基于FastAPI \citep{fastapi} 实现，负责航迹与XML配置文件的生成，并通过子进程调用HELIOS++引擎执行仿真，经WebSocket \citep{WebSocket} 将运行日志实时推送到前端。用户通过浏览器即可完成从场景构建、航路设计、仿真执行到结果可视化的全流程操作，无需安装任何客户端软件。

在项目开发过程中，我们全程采用 Git 与 GitHub 作为版本控制与协作平台。所有代码变更均通过分支（branch）提交，并经由 Pull Request（PR）发起合并请求，由仓库唯一的 CODEOWNER（王孜悠，WZYivan）进行代码审查（review）后方可合并至主分支。这一机制借助 GitHub 的分支保护规则与 CODEOWNERS 文件实现，确保了每次合入的代码都经过审核，历史记录清晰可溯。同时，团队成员通过 Fork + PR 的工作流贡献代码，既保证了主仓库的稳定性，也降低了协作门槛。CODEOWNER 不仅负责代码审查与合并，还承担了系统架构设计、技术路线决策、核心代码实现（包括三维场景渲染、航点交互、航迹生成、HELIOS++集成、点云解析等关键模块）、文档编写与项目管理统筹等角色。从技术选型（Vue 3 + Three.js + FastAPI，摒弃QGIS路线）到开发规范（Git工作流、分支策略、CHANGELOG与TODO维护），均由CODEOWNER主导推进。这种本地开发与云端协作相结合的现代化软件开发实践，不仅提升了团队的协作效率，也充分体现了软件工程最佳规范在本项目中的落地。详细的贡献可见\refApp{app:git}或直接访问\FeHALSRepo。

本文的主要贡献包括：（1）提出一种面向HELIOS++的Web端三维航路规划交互方案，通过鼠标点击与拖拽即可在三维场景中完成航点设定与航线编辑；（2）实现仿真参数的图形化配置与XML、航迹文件的自动生成，屏蔽了HELIOS++底层配置文件的复杂性；（3）实现仿真执行过程的日志实时推送与点云结果的可视化渲染，形成完整的闭环工作流。

本文其余部分组织如下：\refSec{sec:tech}介绍相关技术；\refSec{sec:design}进行系统需求分析与总体设计；第\refSec{sec:impl}阐述系统各模块的设计与实现；\refSec{sec:conclusion}给出结论与展望。

\begin{figure}[pos=!htbp]
    \centering
    \includegraphics[width=\textwidth]{figs/fehals_overview.jpg}
    \caption{FeHALS的Web前端与仿真结果}
    \label{fig:fehals_overview}
\end{figure}

\section{相关技术}\label{sec:tech}
\subsection{HELIOS++仿真框架}
HELIOS++是一款基于C++的开源激光雷达仿真框架，采用光线追踪技术对激光脉冲在三维场景中的传播过程进行建模，能够模拟机载、地面、车载与无人机等多种平台，支持全波形（full-waveform）输出与LAS、LAZ、XYZ等多种点云格式。其仿真配置以XML文件组织：场景（scene）文件定义参与扫描的三维地物模型；平台（platform）与扫描器（scanner）目录文件定义飞行平台与扫描仪的物理参数；扫描任务（survey）文件将场景、平台、扫描器与航迹（trajectory）关联起来，并设定脉冲频率、扫描频率与扫描角等核心参数。HELIOS++提供命令行接口，通过指定survey文件、资源目录（--assets）与输出目录（--output）即可启动仿真。由于配置文件结构复杂且缺乏图形化交互，直接使用HELIOS++进行航路规划与参数调试对非专业用户而言存在较大困难。HELIOS++官方团队已开发了基于QGIS的图形化插件AEOS，将HELIOS++的仿真配置与执行集成到QGIS桌面环境中，为用户提供了GIS工作流内的图形化操作界面。本文提出的FeHALS则采用Web架构，在浏览器中提供轻量级的航路规划与仿真交互体验，与AEOS形成互补。

\subsection{Web三维可视化与Three.js}
Three.js是一个基于WebGL \citep{webgl} 的开源JavaScript三维图形库，封装了场景（Scene）、相机（Camera）、渲染器（Renderer）与网格（Mesh）等核心对象，支持OBJ、GLTF、STL等多种三维模型格式的加载，并提供了轨道控制（OrbitControls）、拖拽控制（DragControls）与射线拾取（Raycaster）等交互组件。借助Three.js，可在浏览器中构建轻量级的三维可视化场景，实现模型展示、点云渲染与鼠标交互拾取，为航路规划提供直观的操作界面。本系统前端即基于Three.js实现三维场景渲染与航点交互。

\subsection{FastAPI与后端异步任务}
FastAPI是一个基于Python的高性能Web框架，基于类型注解与Pydantic \citep{pydantic} 实现自动的数据校验与接口文档生成，并原生支持异步编程与WebSocket。本系统后端基于FastAPI构建REST API \citep{fielding2000rest} 与WebSocket日志通道，通过asyncio子进程异步调用HELIOS++可执行文件，实时捕获其标准输出并以WebSocket推送到前端，实现仿真过程的实时监控。

\section{系统需求分析与设计}\label{sec:design}
\subsection{用户需求分析}
本系统的目标用户为航测工程师、遥感研究人员以及相关专业师生。其核心需求可归纳为三个方面：一是直观的航路规划，用户希望在不编写配置文件的情况下，通过三维场景中的鼠标交互即可完成航点设定与航线编辑；二是快速的仿真执行，用户能够通过图形化界面配置平台与扫描仪参数，一键调用HELIOS++引擎并实时查看运行状态；三是便捷的结果可视化，仿真生成的点云能够自动加载到三维场景中，并支持按高度或强度着色以及尺寸、透明度的调节。

\subsection{系统总体架构}
系统采用浏览器/服务器（B/S）架构，分为前端、后端与仿真引擎三层，其总体架构如\refFig{fig:gernal_fehals_arch}所示。前端基于Vue 3构建，负责三维场景渲染、航点交互编辑、参数配置与结果展示，通过HTTP \citep{rfc9110,rfc9113,rfc9114} 与WebSocket与后端通信；后端基于FastAPI，负责接收前端请求，生成HELIOS++所需的航迹文件与XML配置文件，调用HELIOS++可执行文件执行仿真，并解析与返回点云结果；仿真引擎层为HELIOS++，独立于前后端，由后端以子进程方式调用。各模块间通过明确的REST API与WebSocket消息格式进行数据交换，形成“数据导入 $\rightarrow$ 航路规划 $\rightarrow$ 仿真配置 $\rightarrow$ 执行控制 $\rightarrow$ 结果展示 $\rightarrow$ 日志诊断”的工作流，如\refFig{fig:gernal_fehals_workflow}所示。

\begin{figure}[pos=!htbp]
    \centering
    \includegraphics[width=0.85\textwidth]{figs/gernal_fehals_arch.pdf}
    \caption{FeHALS总体架构}
    \label{fig:gernal_fehals_arch}
\end{figure}

\begin{figure}[pos=!htbp]
    \centering
    \includegraphics[width=0.85\textwidth]{figs/gernal_fehals_workflow.pdf}
    \caption{FeHALS数据流}
    \label{fig:gernal_fehals_workflow}
\end{figure}

\subsection{功能优先级划分}
结合开发环境与工程约束，系统功能划分为三个优先级：核心功能（MVP）包括三维场景渲染、模型加载、交互式航点编辑、航迹导出、仿真参数配置、仿真执行、日志控制台与点云渲染；重要功能包括场景要素属性编辑、自动航线生成、航高自动计算、仿真预设方案、点云特征统计与点云下载；拓展功能包括仿真进度实时反馈、点云覆盖度分析、多任务队列管理与仿真过程动画。本文撰写时仅实现了核心功能，其余功能仍在开发中，最新动态请关注\FeHALSRepo。

\section{系统设计与实现}\label{sec:impl}

\subsection{系统模块划分}
系统实现按前后端分层组织，详细的功能划分如\refFig{fig:detailed_fehals_arch}所示。前端基于Vue 3与Node.js \citep{nodejs} 生态构建，目录下包含组件（components）、状态管理（stores）与组合式函数（composables）三部分：组件层包含三维场景Scene3D.vue、参数配置面板ControlPanel.vue、航点列表WaypointList.vue与日志控制台LogConsole.vue；stores层基于Pinia维护场景（scene）、航点（waypoints）与仿真（simulation）三份响应式状态；composables层封装Three.js场景管理（useThreeScene）、航点交互（useWaypoints）与后端接口调用（useHeliosAPI）。后端基于FastAPI构建，api目录下包含REST路由routes.py与WebSocket端点websocket.py，models目录定义Pydantic数据模型schemas.py，services目录封装航迹生成trajectory\_generator.py、配置生成config\_generator.py、HELIOS++调用helios\_service.py与点云解析pointcloud\_parser.py。

\begin{figure}[pos=!htbp]
    \centering
    \includegraphics[width=\textwidth]{figs/detialed_fehals_arch.pdf}
    \caption{FeHALS详细架构}
    \label{fig:detailed_fehals_arch}
\end{figure}

\subsection{三维场景渲染与航点交互}
三维场景基于Three.js构建，Scene3D.vue负责初始化场景、相机、渲染器与灯光，并添加地面网格与拾取平面；useThreeScene.js以单例形式管理渲染循环、模型与点云的增删以及航点的三维表示。场景使用Z-up坐标系（与HELIOS++一致，z为高程轴），通过OrbitControls实现相机的旋转、平移与缩放；模型加载采用OBJLoader、GLTFLoader与STLLoader，加载完成后根据模型包围盒自动调整相机视角，并基于模型顶点包围盒推断其up轴（Y-up模型自动旋转至Z-up）。航点交互通过自绘拖拽实现：鼠标左键点击经Raycaster与地面平面求交，在交点处添加球体表示的航点（恒为z=0贴地），相邻航点之间以白色折线连接；拖拽航点时锁定相机（`controls.enabled=false`），确保不干扰视角操作；通过Delete键或右键菜单删除航点。

\subsection{前端组件与状态管理}
侧边栏采用多Tab布局，包含仿真参数、点云、模型列表、航迹与设置五个Tab。ControlPanel.vue提供仿真参数表单，分为「载体参数」与「传感器参数」两个模块，各参数旁显示有效范围，只读参数（如光束发散角、最小测程）以灰色禁用显示；参数数据来源于从HELIOS++扫描器与平台目录文件中提取的规格库，切换平台类型时自动重置参数至默认值。WaypointList.vue以列表形式展示航点序号与坐标，支持单个删除与整体清空。LogConsole.vue展示分级日志（INFO/WARNING/ERROR），支持自动滚动与清空。ModelList.vue管理已加载模型，支持可见性切换、bbox线框显示（含尺寸标签）与移除。PointCloudPanel.vue提供点云渲染参数调节（尺寸、透明度、着色模式）与下载。SettingsPanel.vue显示各缓存目录的统计信息并支持清理。三个Pinia store分别维护场景模型与点云渲染参数（scene）、航点坐标数组（waypoints）以及仿真任务状态、进度、日志与结果（simulation）。前端通过useHeliosAPI.js中的Axios实例调用后端REST接口，并通过connectLogWS建立WebSocket连接以接收仿真日志。

\subsection{后端服务与接口}
后端基于FastAPI实现，routes.py提供模型（上传/列表/删除）、航迹（生成/列表/下载）、配置（生成）、仿真（执行/状态/日志/取消）与结果（元数据/下载）五类REST接口，请求与响应由schemas.py中的Pydantic模型校验；上传的模型与生成的航迹、配置、结果统一存放于静态目录并由FastAPI静态文件服务对外提供访问。websocket.py实现/ws/logs/\{task\_id\}端点，连接后先回放已有日志，任务运行期间持续推送增量日志与进度消息。main.py完成CORS中间件、静态目录挂载与路由装配。

\subsection{航迹与配置文件生成}
trajectory\_generator.py将航点数组按添加顺序、以1秒时间间隔生成HELIOS++原生航迹文件（.trj，CSV格式，列顺序为时间、roll、pitch、yaw、x、y、z），所有航点统一为恒定飞行高度（仅用x、y定义水平路径）。系统还支持弓字形自动航迹生成：用户在三维场景中点击两个角点定义矩形区域，系统根据设定航线间距自动生成往返覆盖的全覆蓋航线。config\_generator.py负责生成HELIOS++所需的两个XML文件：场景文件通过objloader滤镜引用三维模型，并指定模型的up轴（Y-up自动旋转至Z-up）；扫描任务（survey）文件引用场景、平台与扫描器目录，并将平台类型映射为对应的平台目录项（UAV与Airborne均使用`copter\_linearpath`平台、`riegl\_vux-1uav`扫描器），将脉冲频率（kHz）与扫描角（±半角）换算为HELIOS++所需的脉冲频率（Hz）与总扫描角。

\subsection{HELIOS++调用与点云渲染}
helios\_service.py维护仿真任务注册表，通过asyncio.create\_subprocess\_exec以子进程方式调用helios++可执行文件，指定资源目录（--assets）、输出目录（--output）与输出格式参数，实时捕获其标准输出/错误流并解析进度百分比，经WebSocket推送至前端，同时支持超时终止与任务取消。仿真启动前，前端对全部参数进行范围校验（数据来源于从HELIOS++扫描器目录文件中提取的参数规格库），超出有效范围时拒绝执行并给出提示。仿真结束后，pointcloud\_parser.py利用laspy \citep{laspy} 解析LAS/LAZ点云或按文本方式解析XYZ点云，提取坐标与强度、计算空间范围，并对超过百万级的点云进行降采样。前端将返回的点云坐标构建为BufferGeometry与Points对象，支持按高度（蓝色至红色渐变）、按强度或固定颜色着色，并可实时调节点尺寸与透明度 \citep{pointCloudRendering}。需要说明的是，仅OBJ格式模型参与HELIOS++仿真，GLTF/STL仅用于三维展示；由于本机HELIOS++构建的LAS输出不可用，系统默认采用XYZ输出格式，详细的数据流见\refFig{fig:detailed_fehals_workflow}。

\begin{figure}[pos=!htbp]
    \centering
    \includegraphics[width=0.85\textwidth]{figs/detailed_fehals_workflow.pdf}
    \caption{FeHALS详细数据流}
    \label{fig:detailed_fehals_workflow}
\end{figure}

\section{结论与展望}\label{sec:conclusion}
本文针对HELIOS++原生命令行交互复杂、缺乏三维可视化航路规划能力的问题，设计并实现了一个基于Web的三维可视化航路规划与激光仿真系统FeHALS。系统采用Vue 3 + Three.js前端与FastAPI后端，实现了三维场景渲染、模型加载、交互式航点编辑、弓字形自动航迹生成、航迹文件导出、仿真参数图形化配置（含参数规格库与范围校验）、HELIOS++引擎调用、日志实时推送、点云渲染与导出、模型列表管理（可见性/bbox/移除）、缓存清理等核心功能，通过典型场景验证了系统的可行性与实用性，有效降低了HELIOS++的使用门槛。

在项目管理方面，本文作者作为仓库唯一的CODEOWNER，统筹了10名成员的协作开发，累计完成40余次提交与8次Pull Request合并。从技术选型决策（摒弃QGIS路线转向Web架构）、开发规范制定（Git工作流、分支策略、CODEOWNERS保护规则）、核心代码实现（三维场景、航点交互、航迹生成、HELIOS++集成、点云解析等模块），到文档编写、团队进度管理、代码审查与质量把控，均由作者主导推进。各成员的具体贡献类型详见首页的贡献声明与\refApp{app:git}，详细的提交历史与代码变更可访问\FeHALSRepo。

后续工作将在以下方面展开：（1）实现HELIOS++运行环境自动检测与诊断，系统启动时检测helios++可执行文件是否可用、资源目录（--assets）是否完整，并在前端设置面板中直观显示检测结果，缺失时提供从源码构建到配置的完整引导提示；（2）提供应用环境一键部署方案，包括Conda环境配置文件（environment.yml）与自动安装脚本，实现Python依赖、Node.js依赖与前端构建的自动化，并可扩展Docker容器化部署方式；（3）引入航高自动计算，根据场景中模型最高点自动推荐安全飞行高度；（4）实现仿真任务队列管理，支持多任务顺序执行与并发调度，显示任务列表与状态；（5）增强点云后处理能力，提供覆盖度分析、密度热力图与特征统计等分析功能；（6）优化大规模点云的渲染性能，采用分层细节（LOD）与八叉树等结构支撑千万级点云的流畅交互；（7）拓展扫描器与平台规格库，支持用户自定义参数配置；（8）实现仿真过程动画，逐点或逐线显示点云生成过程。

\printcredits
\newpage
\appendix

\section{详细的贡献}\label{app:git}

\begin{figure}[pos=!htbp]
    \centering
    \includegraphics[width=\textwidth]{figs/gitgraph.pdf}
    \caption{基于git历史的详细贡献图，用户名的对应关系见首页脚注}
\end{figure}

\FloatBarrier

\section{FeHALS示例工作流}

\subsection{模型导入}
\subsection{仿真参数}
\subsection{航迹绘制}
\subsection{仿真执行}
\subsection{点云导出}

\FloatBarrier
\newpage

\bibliographystyle{cas-model2-names}
\bibliography{cas-refs}

\end{document}
